% GENERATED FILE. Edit canonical lessons, then run npm run generate:books.
% canonical-source-hash: fnv1a64:f15eb285d6144b04
% canonical-lessons: UR-C302-afvah, UR-C302-charcha, UR-C302-dahshat, UR-C302-bharkana, UR-C302-avam

\chapter{A Rumour Spreads}
\label{ch:ur-a-rumour-spreads}
\hlchaptermodality{302}

\begin{chapteropening}
\textbf{By the end of this chapter:} \emph{I can say how a rumour spread and how people reacted.}

Complete the last lesson of chapter 302: I can say how a rumour spread and how people reacted.
\end{chapteropening}

\section[\texorpdfstring{\ur{افواہ} (afvāh)}{افواہ (afvāh)}]{\ur{افواہ} (afvāh) — a rumour}

\label{lesson:UR-C302-afvah}



\begin{quote}
Before the new one: say the Urdu for a detail, then the Urdu for to have written down.
\end{quote}

\subsection*{You'll want to know: \ur{افواہ}}
\textbf{\ur{افواہ}} — \emph{afvāh} — \textquotedblleft{}a rumour\textquotedblright{}.

\begin{quote}
\ur{شہر} \ur{میں} \ur{ایک} \ur{افواہ} \ur{تھی۔}
\end{quote}

\begin{quote}
\emph{shahr meṅ ek afvāh thī.}
\end{quote}

\begin{quote}
There was a rumour going round the city.
\end{quote}

Feminine.

\begin{cousinweb}
\textbf{\ur{افواہ}} (\emph{afvāh}) is usually traced, through Persian, to Arabic \emph{afwāh}, \textquotedblleft{}mouths\textquotedblright{}: what passes from mouth to mouth. It is feminine.
\end{cousinweb}

\begin{grammarlens}[title={What you've built}]
One more word for telling what happened.
\end{grammarlens}

\subsection*{Your turn}
\begin{itemize}
\raggedright
  \item \emph{Say it:} \emph{afvāh}
  \item \emph{Say it:} \emph{afvāh}, once more
  \item \emph{Say it:} \emph{shahr meṅ ek afvāh thī.} — There was a rumour going round the city.
  \item \emph{Recall:} say the Urdu for a detail, then the Urdu for to have written down, then say \emph{afvāh} again
\end{itemize}

\begin{tcolorbox}[breakable,colback=teal!4,colframe=teal!35!black,title={Before you move on}]
What does \ur{افواہ} mean? (\textquotedblleft{}A rumour\textquotedblright{}.) Say it once more.
\end{tcolorbox}
\section[\texorpdfstring{\ur{چرچا} (charchā)}{چرچا (charchā)}]{\ur{چرچا} (charchā) — talk of the town; a much-discussed matter}

\label{lesson:UR-C302-charcha}



\begin{quote}
Before the new one: say the Urdu for to have written down, then the Urdu for a rumour.
\end{quote}

\subsection*{You'll want to know: \ur{چرچا}}
\textbf{\ur{چرچا}} — \emph{charchā} — \textquotedblleft{}talk of the town; a much-discussed matter\textquotedblright{}.

\begin{quote}
\ur{ہر} \ur{جگہ} \ur{اس} \ur{بات} \ur{کا} \ur{چرچا} \ur{تھا۔}
\end{quote}

\begin{quote}
\emph{har jagah is bāt kā charchā thā.}
\end{quote}

\begin{quote}
Everyone everywhere was talking about it.
\end{quote}

Masculine.

\begin{cousinweb}
\textbf{\ur{چرچا}} (\emph{charchā}) is usually traced to Sanskrit \emph{carcā}, \textquotedblleft{}repetition, a going over\textquotedblright{}, in Monier-Williams; in Urdu it is the talk that goes round. It is masculine: \textbf{\ur{چرچا} \ur{تھا}}.
\end{cousinweb}

\begin{grammarlens}[title={What you've built}]
One more word for telling what happened.
\end{grammarlens}

\subsection*{Your turn}
\begin{itemize}
\raggedright
  \item \emph{Say it:} \emph{charchā}
  \item \emph{Say it:} \emph{charchā}, once more
  \item \emph{Say it:} \emph{har jagah is bāt kā charchā thā.} — Everyone everywhere was talking about it.
  \item \emph{Recall:} say the Urdu for to have written down, then the Urdu for a rumour, then say \emph{charchā} again
\end{itemize}

\begin{tcolorbox}[breakable,colback=teal!4,colframe=teal!35!black,title={Before you move on}]
What does \ur{چرچا} mean? (\textquotedblleft{}Talk of the town; a much-discussed matter\textquotedblright{}.) Say it once more.
\end{tcolorbox}
\section[\texorpdfstring{\ur{دہشت} (dahshat)}{دہشت (dahshat)}]{\ur{دہشت} (dahshat) — terror; panic}

\label{lesson:UR-C302-dahshat}



\begin{quote}
Before the new one: say the Urdu for a rumour, then the Urdu for talk of the town.
\end{quote}

\subsection*{You'll want to know: \ur{دہشت}}
\textbf{\ur{دہشت}} — \emph{dahshat} — \textquotedblleft{}terror; panic\textquotedblright{}.

\begin{quote}
\ur{لوگوں} \ur{میں} \ur{دہشت} \ur{تھی۔}
\end{quote}

\begin{quote}
\emph{logoṅ meṅ dahshat thī.}
\end{quote}

\begin{quote}
People were terrified.
\end{quote}

Feminine.

\begin{cousinweb}
\textbf{\ur{دہشت}} (\emph{dahshat}) is usually traced, through Persian, to Arabic \emph{dahsha}, \textquotedblleft{}bewilderment, alarm\textquotedblright{}. It is a fear that grips a whole crowd.
\end{cousinweb}

\begin{grammarlens}[title={What you've built}]
One more word for telling what happened.
\end{grammarlens}

\subsection*{Your turn}
\begin{itemize}
\raggedright
  \item \emph{Say it:} \emph{dahshat}
  \item \emph{Say it:} \emph{dahshat}, once more
  \item \emph{Say it:} \emph{logoṅ meṅ dahshat thī.} — People were terrified.
  \item \emph{Recall:} say the Urdu for a rumour, then the Urdu for talk of the town, then say \emph{dahshat} again
\end{itemize}

\begin{tcolorbox}[breakable,colback=teal!4,colframe=teal!35!black,title={Before you move on}]
What does \ur{دہشت} mean? (\textquotedblleft{}Terror; panic\textquotedblright{}.) Say it once more.
\end{tcolorbox}
\section[\texorpdfstring{\ur{بھڑکانا} (bhaṛkānā)}{بھڑکانا (bhaṛkānā)}]{\ur{بھڑکانا} (bhaṛkānā) — to stir up; to incite}

\label{lesson:UR-C302-bharkana}



\begin{quote}
Before the new one: say the Urdu for talk of the town, then the Urdu for terror.
\end{quote}

\subsection*{You'll want to know: \ur{بھڑکانا}}
\textbf{\ur{بھڑکانا}} — \emph{bhaṛkānā} — \textquotedblleft{}to stir up; to incite\textquotedblright{}.

\begin{quote}
\ur{افواہ} \ur{نے} \ur{لوگوں} \ur{کو} \ur{بھڑکایا۔}
\end{quote}

\begin{quote}
\emph{afvāh ne logoṅ ko bhaṛkāyā.}
\end{quote}

\begin{quote}
The rumour stirred people up.
\end{quote}

Past, with \ur{نے} (he or she …): \ur{اس} \ur{نے} \ur{لوگوں} \ur{کو} \ur{بھڑکایا}, \emph{us ne logoṅ ko bhaṛkāyā}.

\begin{cousinweb}
\textbf{\ur{بھڑکانا}} (\emph{bhaṛkānā}) is the causing partner of \textbf{\ur{بھڑکنا}}, to flare up, which CDIAL lists for Hindi under \emph{*bhaṭ} (CDIAL 9365): to make a fire, or a crowd, flare.
\end{cousinweb}

\begin{grammarlens}[title={What you've built}]
One more word for telling what happened.
\end{grammarlens}

\subsection*{Your turn}
\begin{itemize}
\raggedright
  \item \emph{Say it:} \emph{bhaṛkānā}
  \item \emph{Say it:} \emph{bhaṛkānā}, once more
  \item \emph{Say it:} \emph{afvāh ne logoṅ ko bhaṛkāyā.} — The rumour stirred people up.
  \item \emph{Recall:} say the Urdu for talk of the town, then the Urdu for terror, then say \emph{bhaṛkānā} again
\end{itemize}

\begin{tcolorbox}[breakable,colback=teal!4,colframe=teal!35!black,title={Before you move on}]
What does \ur{بھڑکانا} mean? (\textquotedblleft{}To stir up; to incite\textquotedblright{}.) Say it once more.
\end{tcolorbox}
\section[\texorpdfstring{\ur{عوام} (avām)}{عوام (avām)}]{\ur{عوام} (avām) — the public; the common people}

\label{lesson:UR-C302-avam}



\begin{quote}
Before the new one: say the Urdu for terror, then the Urdu for to stir up.
\end{quote}

\subsection*{You'll want to know: \ur{عوام}}
\textbf{\ur{عوام}} — \emph{avām} — \textquotedblleft{}the public; the common people\textquotedblright{}.

\begin{quote}
\ur{عوام} \ur{پریشان} \ur{تھے۔}
\end{quote}

\begin{quote}
\emph{avām pareshān the.}
\end{quote}

\begin{quote}
The public were worried.
\end{quote}

Masculine, used in the plural.

\begin{cousinweb}
\textbf{\ur{عوام}} (\emph{avām}) is usually traced, through Persian, to Arabic \emph{ʿawāmm}, \textquotedblleft{}the common people\textquotedblright{}. It takes a plural verb: \textbf{\ur{عوام} \ur{پریشان} \ur{تھے}}.
\end{cousinweb}

\begin{grammarlens}[title={What you've built}]
One more word for telling what happened.
\end{grammarlens}

\subsection*{Your turn}
\begin{itemize}
\raggedright
  \item \emph{Say it:} \emph{avām}
  \item \emph{Say it:} \emph{avām}, once more
  \item \emph{Say it:} \emph{avām pareshān the.} — The public were worried.
  \item \emph{Recall:} say the Urdu for terror, then the Urdu for to stir up, then say \emph{avām} again
\end{itemize}

\begin{tcolorbox}[breakable,colback=teal!4,colframe=teal!35!black,title={Before you move on}]
What does \ur{عوام} mean? (\textquotedblleft{}The public; the common people\textquotedblright{}.) Say it once more.
\end{tcolorbox}
